% TOP_PROBLEM: {"id": 30006678, "title": "NP versus P/poly", "category_id": 15, "category": {"id": 15, "name": "computer_science", "display_name": "Computer Science", "description": "Computational complexity, algorithms, and theoretical CS.", "slug": "computer-science", "order_index": 15, "created_at": "2026-07-31T15:26:25.671Z"}, "status": "open"}
% FIRST_POSTED: 2026-09-27
% !TeX program = lualatex
% ATTEMPT_STATUS: unresolved
\documentclass[11pt]{article}
\usepackage[margin=25mm]{geometry}
\usepackage{fontspec}
\setmainfont{Latin Modern Roman}
\usepackage{amsmath,amssymb,mathtools}
\usepackage{unicode-math}
\setmathfont{Latin Modern Math}
\usepackage{xurl,hyperref}
\hypersetup{colorlinks=true,urlcolor=blue}
\setcounter{secnumdepth}{0}
\setlength{\parindent}{0pt}
\setlength{\parskip}{5pt}
\setlength{\emergencystretch}{3em}
\begin{document}
\section*{\texorpdfstring{NP versus P/poly}{NP versus P/poly}}\label{30006678}
\subsection{Definitions and mathematical statement}
A Boolean circuit is a finite directed acyclic computation with binary AND/OR gates, unary NOT gates, input bits, constants, and one output. Define $\mathsf{P/poly}$ by
\[
 L\in\mathsf{P/poly}\ \Longleftrightarrow\ \exists C,k\ \forall n\ge0\ \exists C_n:
 |C_n|\le C(n+1)^k,\quad C_n(x)=1_L(x)\ (|x|=n).
\]
All nodes are counted. No algorithm for generating $C_n$ is required. With $\mathsf{NP}$ defined by polynomially checkable witnesses, determine whether $\mathsf{NP}\subseteq\mathsf{P/poly}$. Different languages may have different size bounds.
\par\smallskip\textit{Formulation and concept references:} \href{https://theory.cs.princeton.edu/complexity/book.pdf}{[S1]}.

\subsection{Short English statement}
Can every efficiently verifiable decision problem be solved by small circuits when a different circuit may be chosen separately for each input length?
\subsection{Research attempt}
\paragraph{Scope and literature check.}
This unresolved attempt by GPT-6 Astra Ultra was prepared on 27 September
2026 by checking the catalogue, deriving a counting construction, and then
looking for the first place where it fails to yield an NP lower bound.
The circuit basis and the convention that input and constant nodes count
are retained. Unbounded search for the appropriate circuit is permitted in
$\mathsf{P/poly}$; a lower bound on the time needed to \emph{find} circuits
would not settle this question.

The circuit/advice equivalence and Karp--Lipton implication in [S1],
Sections 6.1--6.2, were checked. The latter says that
$\mathsf{NP}\subseteq\mathsf{P/poly}$ collapses the polynomial hierarchy to
its second level; a collapse is a consequence, not a contradiction already
known to be impossible. The original nonuniform-to-uniform work is [E1].
The counting argument below is deliberately elementary and self-contained;
it is not claimed to improve known circuit lower bounds.

Write $s_L(n)$ for the least number of nodes in a circuit for the length-$n$
characteristic function of $L$. A negative answer needs
\[
 \exists L\in\mathsf{NP}\ \forall C>0\ \forall k\ge1\ \exists n:
                       s_L(n)>C(n+1)^k.
\]
Equivalently, one fixed NP language must beat each fixed polynomial bound
on infinitely many lengths. The word ``fixed'' applies to the language and
its verifier before choosing the exponent to be defeated. A family of
languages indexed by that exponent does not suffice.

\paragraph{Attack plan: make the counting witness explicit.}
There are $2^{2^n}$ Boolean functions on $n$ bits but only
$2^{O(s\log(n+s))}$ circuits with at most $s$ gates. This proves that most
functions require large circuits, but it does not place their truth tables
inside a language with an NP verifier. A natural repair is to select the
first missing truth table. The following construction makes that selection
fully uniform, bounds its working storage, and identifies an additional
failure: the selected language still has polynomial-size circuits.

\paragraph{Lemma 12.1: an explicit fixed-exponent diagonalization.}
For every fixed integer $k\ge1$, there is a language $L_k$ in
$\mathsf{PSPACE}\cap\mathsf{P/poly}$ such that, for all sufficiently large
$n$, no circuit of at most $n^k$ nodes computes its length-$n$ restriction.
A deterministic decider can use $O(n^k\log(n+2))$ work bits and
$2^{O(n^k\log(n+2))}$ time.

\emph{Proof.} Put $s=n^k$ and $q=n+s+2$. Every circuit with at most $s$
internal gates can be listed in topological order. For each gate allow at
most four types and two source indices in a pool of size $q$; unused source
slots and invalid choices only increase the count. There are at most $q$
choices for the output. Summing over the possible number of gates yields
the explicit upper bound
\begin{equation}\label{a12:count}
                 N=(s+1)q(4q^2)^s.
\end{equation}
Every circuit of at most $s$ \emph{total} nodes occurs in this larger class.
Repeated encodings of the same circuit or function do not hurt the bound.
The integer $N$ has
\[
                    B=\lceil\log_2(N+1)\rceil
                       =O(n^k\log(n+2))
\]
bits. For all sufficiently large $n$, $B<2^n$.

For each $j\in\{0,\ldots,N\}$, let $T_j$ be the $2^n$-bit binary expansion
of $j$, padded by leading zeroes, regarded as a truth table in lexicographic
input order. These $N+1$ tables are distinct. At most $N$ of them can be
computed by the circuits being enumerated. Let $j_n$ be the smallest index
whose table is not computed by any such circuit, and define
\[
                     1_{L_k}(x)=T_{j_n}[x]
                    \qquad(|x|=n).
\]
Use the zero function at the finitely many lengths where $B\ge2^n$.
The choice exists by counting and immediately gives the lower bound.

To compute $j_n$, enumerate $j=0,\ldots,N$. For each $j$, enumerate the
valid circuit descriptions and test whether each disagrees with $T_j$ on
some input. Test equality by traversing all $2^n$ inputs, evaluating the
circuit and one requested bit of $T_j$. Neither a circuit truth table nor
the padded string $T_j$ is stored. The nonzero part of $T_j$ is contained
in its last $B$ positions, so a requested bit is obtained from the $B$-bit
integer $j$ and an $n$-bit position counter. A circuit description uses
$O(s\log q)$ bits; its evaluation uses polynomially many additional bits,
in fact $O(s+n)$ for gate values and the current input. Reuse this storage
for the successive candidates. The circuit enumeration can likewise be
implemented with $O(s\log q)$ counters and descriptions.

There are at most $N+1$ candidates, at most $N$ circuit descriptions,
and $2^n$ input evaluations per comparison. Thus the time is bounded by
$(N+1)N2^n\operatorname{poly}(n,s)$, which is
$2^{O(n^k\log(n+2))}$, and the work space is $O(B)$. Both exponents refer
to the previously fixed integer $k$.

Finally, $T_{j_n}$ contains at most $B$ ones. A circuit can hardwire all
accepted strings of length $n$, test equality with each using $O(n)$ gates,
and OR the results. Shared negated inputs and binary trees implement this
with $O(nB+n)$ total nodes. Therefore
\begin{equation}\label{a12:upper}
                s_{L_k}(n)=O(n^{k+1}\log(n+2)),
\end{equation}
up to adjustment at finitely many lengths. This is a polynomial bound for
fixed $k$, so $L_k\in\mathsf{P/poly}$. The enumeration algorithm gives
$L_k\in\mathsf{PSPACE}$.\hfill$\square$

This is more informative than saying that counting is nonconstructive.
The counting witness here \emph{is} uniformly computable with a polynomial
space bound, and nevertheless it provably fails the required nonuniform
separation. The construction forces a bound larger than $n^k$, while its
choice of small integer truth tables guarantees a somewhat larger polynomial
upper bound. Both bounds hold for the same language without contradiction.

\paragraph{Where an NP-certificate argument stalls.}
For a candidate $j$ of polynomial bit length, its missing-table property is
\begin{equation}\label{a12:quantifiers}
 \forall\text{ circuits }C\text{ in the size bound}\quad
 \exists x\in\{0,1\}^n:\ C(x)\ne T_j[x].
\end{equation}
Each individual disagreement is efficiently checked. But a list with one
disagreement for every circuit can have exponential length. Guessing just
$j$ does not verify the universal circuit condition. Guessing one circuit
and one disagreement does not verify it either, because that reverses the
quantifier pattern. Selecting the \emph{least} missing table additionally
requires establishing that earlier candidates were not missing.

Equation \eqref{a12:quantifiers} is a precise gap in this route, rather than
a theorem that no NP verifier for $L_k$ could exist. The attempted verifier
has not supplied a polynomial certificate for that universal condition.
Even supplying one would still leave the upper bound \eqref{a12:upper} and
the need to defeat every polynomial for a single language.

\paragraph{Testing two repairs.}
First, let $k$ grow slowly with $n$ and repeat the diagonalization. This can
produce a single computable language whose slices eventually beat each
fixed power. However its stored circuit descriptions have length
$n^{k(n)}$ up to logarithmic factors. If $k(n)$ is unbounded, this is not
bounded by any fixed polynomial. The particular decider above then has no
polynomial-space guarantee, much less a polynomial NP witness bound.
Scheduling a larger exponent at a sparse sequence of lengths does not
change this worst-case requirement.

Second, pad an input of length $n$ to length $N=n^d$, with $d$ fixed, hoping
to pay for the search. A bound larger than $n^k$ then becomes only a bound
larger than $N^{k/d}$, modulo the gates needed to hardwire padding. It still
defeats a fixed exponent. If the padding is large enough to make an
exponential search polynomial in the padded length, that same size change
can make the inherited circuit bound much too small. Neither padding nor
renaming the exponent repairs the quantifiers automatically.

A positive answer would require polynomial circuits for an NP-complete
language, such as SAT, with arbitrary dependence on input length allowed.
No such circuit construction appears here. The absence of an efficient
uniform circuit generator cannot be used as evidence against that answer.

\paragraph{Verification.}
The counting bound deliberately counts a larger gate model than the target,
so the lower bound survives the convention that every node counts. The
upper bound explicitly includes the $n$ input nodes, shared negations,
equality tests, and output tree. Finitely many exceptional lengths can be
hardwired into either description without affecting class membership.

The companion script checks the finite core independently for $n=1,2,3$
with one internal gate: it enumerates every permitted gate and source pair,
forms the distinct truth tables, and finds the least omitted table when one
exists. At one input every function already occurs; at larger inputs the
missing-table selection is well defined. This checks the combinatorial
selection, not the asymptotic resource analysis. In particular no experiment
is presented as evidence that SAT needs large circuits.

\paragraph{Outcome and next targets.}
\textbf{UNRESOLVED.} A fully specified fixed-exponent diagonalization has been
proved, together with a polynomial circuit upper bound for its own output.
This decisively rejects that construction as a separation from
$\mathsf{P/poly}$. Its mathematics is elementary; historical novelty is not
claimed. The remaining target is one NP language with superpolynomial
circuit complexity. Useful next steps would require a property specific to
NP witness relations, a polynomial verifier surviving all exponent choices,
and a lower-bound measure controlling unrestricted circuits rather than
only one syntactic representation. None has been supplied here.

\subsection{Sources}
\begin{itemize}
\item[S1]Arora and Barak, Computational Complexity: A Modern Approach, web draft January8,2007, Definitions6.1–6.3 and Theorem6.13; Aaronson, P ?= NP, Conjecture25.
\url{https://theory.cs.princeton.edu/complexity/book.pdf}.
\textit{Formulation source.}
\item[S2]Boolean Circuits and P/poly, Cornell CS 6810, February 26, 2026.
\url{https://www.cs.cornell.edu/courses/cs6810/2026sp/lec10.pdf}.
\textit{Further reference; not a proof endorsement.}
\item[E1]Richard M. Karp and Richard J. Lipton,
\emph{Some Connections between Nonuniform and Uniform Complexity Classes},
Proceedings of STOC 1980, pp. 302--309.
\url{https://doi.org/10.1145/800141.804678}.
\textit{Original nonuniform-to-uniform reference. The precise second-level
collapse used here was checked in [S1], Theorem 6.13, which credits
improvements by Sipser; consulted 27 September 2026.}
\end{itemize}
\end{document}
